\chapter{\MakeUppercase{Introduction}}

\section{Background}

People turn to the internet when they have medical questions. That is not going to change. What changes is whether the information they find is grounded in verified sources or assembled from statistical patterns that happen to sound authoritative. Most consumer health websites fall into the second category --- written for the median reader, not the specific person asking, and updated infrequently if at all.

\ac{ai}-powered \ac{qa} systems have existed in various forms since the 1970s. MYCIN~\cite{shortliffe1974mycin}, developed at Stanford, used hand-coded rules to recommend antibiotic treatments for blood infections. Within its narrow domain it performed well, but maintaining the rule base required continuous expert input and the system could not generalise beyond what its authors had explicitly encoded. That same limitation --- that medical knowledge is too large and changes too quickly to encode manually --- is one reason this project leans on retrieval instead of asking the generator to hold all clinical detail in its weights.

Modern \ac{llm} systems solve the knowledge encoding problem but introduce a different one. A model trained on web-scale text has encountered a large volume of medical content and can produce fluent, specific-sounding answers. The difficulty is that it cannot reliably distinguish between what it knows accurately and what it has reconstructed from training distributions. \ac{rag}~\cite{lewis2020rag}, as formulated by Lewis et al.\ in 2020, avoids memory-only answers by supplying retrieved documents at query time and conditioning generation on that text. Figure \ref{fig:rag_pipeline} sketches the two-stage flow.

\begin{figure}[!ht]
    \centering
    \includegraphics[width=0.95\linewidth,height=0.85\textheight,keepaspectratio]{images/rag_pipeline_final.png}
    \caption{Retrieval pulls passages from the knowledge base; the language model then generates an answer using that text as context.}
    \label{fig:rag_pipeline}
\end{figure}

A practical constraint runs through every design decision in this project: running a 7B-parameter model on a remote server costs money and routes queries through a third party. For medical questions, sending private data over the wire is a genuine privacy concern. You don't want your symptoms ending up in an external log. TinyLlama 1.1B~\cite{zhang2024tinyllama} is small enough to run on \ac{cpu} hardware that most students and small clinics already own. This system uses TinyLlama as its default model, with BioMistral-7B~\cite{labrak2024biomistral} available as an alternative where the latency tradeoff is acceptable.

On the retrieval side, hybrid search combining sparse (\ac{bm25})~\cite{robertson2009bm25} and dense~\cite{reimers2019sbert} methods consistently outperforms either method alone on open-domain \ac{qa} benchmarks~\cite{ma2022hybrid}. The fusion method used here is Reciprocal Rank Fusion (\ac{rrf})~\cite{cormack2009rrf}, which merges ranked lists without requiring score normalisation across retrieval systems.

A \ac{qa} system that gives no indication of its own uncertainty isn't just unhelpful; in a medical context, it's potentially harmful. A confidently wrong answer can easily mislead a patient. Multi-signal confidence estimation~\cite{guo2017calibration} is included so the stack can flag when it ought to be less certain. The five signals used here cover retrieval quality, generation consistency, source agreement, \ac{nli}-based hallucination detection~\cite{he2021deberta}, and medical entity coverage. Platt scaling~\cite{platt1999probabilistic} maps the combined score to a calibrated display value.

\section{Motivation}

The motivation for this project started with a straightforward observation: when patients ask general search engines about their own health conditions, they rarely get answers tailored to their situation. Results for a query like ``knee swelling after a fall'' are written for a broad audience, padded with disclaimers, and stripped of specificity. A retrieval system indexed on real patient-doctor conversation data has a reasonable chance of matching a specific question pattern to a previously answered case with similar clinical context.

There is also an academic angle. \ac{rag} for general domains is well studied, but healthcare-specific \ac{rag} with explicit confidence signals and hallucination detection is less common. Most published medical \ac{qa} systems report accuracy on structured benchmarks such as MedQA~\cite{jin2021medqa} \allowbreak{}or PubMedQA~\cite{jin2019pubmedqa} without reporting calibration or uncertainty estimates. A system that says ``80\% confident'' is more useful than one that says ``here is the answer'' --- provided that 80\% is actually calibrated to observed accuracy.

The technical interest centred on running this entire stack on \ac{cpu}. Most published \ac{rag} systems assume \ac{gpu} availability for retrieval and generation alike. Making \ac{bm25} and dense retrieval and \ac{llm} generation work at acceptable quality on a machine without a \ac{gpu} required specific choices: \ac{bm25} caching via pickle file, \ac{lru} caching for embedding lookups, and \ac{gguf} quantisation for the 7B model.

\section{Scope of the project}

\subsection{Functional scope}

The system accepts free-text medical questions in English and returns an answer grounded in retrieved documents, with source attribution, a confidence score, and a standard safety disclaimer. Four query types are handled: factual definitions, symptom-based questions, research questions drawn from \allowbreak{}PubMedQA abstracts, and clinical questions from MedMCQA. Questions involving prescription dosages, emergency triage, or content caught by the safety guardrail are redirected to clinical care rather than answered directly.

\subsection{Non-functional scope}

The system runs without a \ac{gpu} and without internet access at query time. User queries are not logged or stored. Response latency averages 106 seconds on a warm path on the test hardware, which places it outside the range of real-time consumer applications. This is a research prototype; latency and output quality would both improve substantially on \ac{gpu} hardware or with a more specialised model.

\subsection{Technical scope}

I built the system to handle the full pipeline from knowledge base construction and hybrid retrieval to generation, \ac{xai} confidence scoring, hallucination detection, a REST \ac{api} layer, and a React frontend. The project does not attempt clinical validation, regulatory compliance, or multi-turn conversation memory. The knowledge base is static once built; there is no mechanism for continuous ingestion of new literature.

\subsection{Assumptions and constraints}

All development and local evaluation ran on \ac{cpu}-only hardware. \ac{gpu} compute was used for two tasks: QLoRA fine-tuning, which ran on a Google Colab \ac{tpu} v5e instance, and BioMistral evaluation, which ran on a Colab T4. All three source datasets are publicly available for research use. No proprietary or patient-identifiable data was used at any point in the project.

\medskip
\noindent\textbf{Report structure.} Chapter~2 surveys related work and states the research gap. Chapter~3 records requirements, feasibility, and the evaluation environment. Chapters~4 and~5 present system design and the methodology used for testing. Chapter~6 documents implementation and deployment. Chapter~7 reports results and discussion; Chapter~8 summarises conclusions and extensions.
